\section{Introduction}

The Nicolai map
transforms interacting supersymmetric theories
to non-interacting ones~\cite{NicolaiMap}.
Supersymmetry is considered to be essential for
the existence of these field transformations in view
of the fact that their Jacobian is exactly
cancelled by the fermion partition function.

In lattice gauge theory,
a natural question to ask is
whether there are field transformations that map
the theory to its strong-coupling limit.
In particular, if there are no matter fields,
one is looking for substitutions
\begin{equation}
  U=\trans(V)
\end{equation}
of the gauge field $U$ in the functional integral
whose Jacobian cancels the gauge-field action.
Similarly to the Nicolai map, this kind of
transformation maps the theory to a solvable one, but
supersymmetry is not required and the Jacobian
plays a different r\^ole.

On a finite lattice, and if the gauge group is compact
and connected, the existence of such
trivializing maps is guaranteed by
a general theorem on volume forms on compact manifolds
(see ref.~\cite{MoserZehnder}, Theorem~1.26, for example).
One may be inclined to assume that these transformations
are too complicated to be of any use.
However, as explained later in this paper,
it is possible to build up trivializing maps
by integrating flows in field space, whose generators satisfy
certain partial differential equations.
The latter are quite tractable and
can, to some extent, be solved analytically in the pure
gauge theory. An application of trivializing
maps, which can then be envisaged, is the
acceleration of lattice QCD simulations.

\begin{figure}[htbp]
  \centering
  \includegraphics[width=6.0cm]{images/cycle.pdf}
  \caption{%
The proposed
simulation algorithm for lattice QCD
updates the gauge field $U$ in three
steps, following the arrows in this diagram, where
the Hamilton function used in the HMC step
has the standard kinetic term and includes the Jacobian
of the field transformation $\mathcal{F}$
(cf.\ subsect.~2.4).
}
  \label{fig:cycle}
\end{figure}

The fact that the efficiency of the
available simulation algorithms is unpredictable
has always been a weakness of numerical lattice QCD.
Already a while ago, empirical studies
of the $\Group$ gauge theory by
Del Debbio, Panagopoulos and Vicari~\cite{DelDebbioTauQ}
showed that the
autocorrelation times of observables related to the topological
charge of the gauge field tend to be large and
appear to grow exponentially with the inverse of the lattice spacing.
Moreover,
Schaefer, Sommer and Virotta~\cite{StefanConference}
recently found that the situation is, in this respect, essentially
unchanged when the sea quarks are included in the simulations.

The rapid slowdown of the simulations at small lattice spacings
may conceivably be overcome
by combining approximately trivializing maps
with the HMC simulation algorithm~\cite{HMC}
(see fig.~1).
Since the transformation moves the theory
closer to the strong-coupling limit, where
the HMC algorithm is known to be highly efficient,
the autocorrelation times are, in general, expected to be
reduced in this way. Evidently, for the combined algorithm
to work out in practice, approximately trivializing maps
must be found which are fairly simple and programmable.
One of the goals of the present paper is thus to provide a
solution to this problem.
